固定 $t$, 考虑空间谱方程 $\phi_x=U\phi$, 其中 $U=-\ii k\sg+Q$. 在 $z\in\Sigma\mathbin{\backslash}\{0,\pm\ii q_0,z_a^\pm\}$ 上定义 Jost 解及修正 Jost 解
\begin{equation}\label{eq:jost-normalization}
\phi_\pm(x,t,z)e^{-\ii\theta(x,t,z)\sg}\longrightarrow Y(z),\qquad x\to\pm\infty,
\end{equation}
\begin{equation}\label{eq:modified-jost}
\mu_\pm=\phi_\pm e^{-\ii\theta\sg},\qquad \mu_\pm\longrightarrow Y\quad(x\to\pm\infty).
\end{equation}

\paragraph{读式说明: 为什么要去掉右侧指数.}
这里固定 $t$, 所有空间微分均把 $t,z$ 看作常数. 由原文的相位定义,
\[
\theta_x=-\lambda,
\qquad
e^{\ii\theta\sg}
=\begin{pmatrix}e^{\ii\theta}&0\\0&e^{-\ii\theta}\end{pmatrix},
\qquad
e^{-\ii\theta\sg}
=\begin{pmatrix}e^{-\ii\theta}&0\\0&e^{\ii\theta}\end{pmatrix}.
\]
因此右乘指数就是分别对两列乘标量:
\[
\begin{aligned}
\mu_{\pm,1}&=e^{-\ii\theta}\phi_{\pm,1},
&\phi_{\pm,1}&=e^{\ii\theta}\mu_{\pm,1},\\
\mu_{\pm,2}&=e^{\ii\theta}\phi_{\pm,2},
&\phi_{\pm,2}&=e^{-\ii\theta}\mu_{\pm,2}.
\end{aligned}
\]
这一步只去掉已知背景振荡. 未去掉指数时, 第一列在空间端点含
$e^{-\ii\lambda x}$, 第二列含 $e^{\ii\lambda x}$;
去掉指数之后, 两列的端态分别成为 $Y_1,Y_2$.
把乘积求导逐项写出, 得
\begin{align*}
(\mu_\pm)_x
&=(\phi_\pm)_xe^{-\ii\theta\sg}
+\phi_\pm(-\ii\theta_x\sg)e^{-\ii\theta\sg}\\
&=U\phi_\pm e^{-\ii\theta\sg}
+\ii\lambda\phi_\pm\sg e^{-\ii\theta\sg}\\
&=U\mu_\pm+\ii\lambda\mu_\pm\sg.
\end{align*}
最后一步允许把 $\sg$ 与指数交换, 因为二者都是对角矩阵;
不允许把 $\sg$ 任意移到 $\mu_\pm$ 左侧.

记 $\Delta Q=Q-Q_\infty$. 以下先在分支点之外构造解析列, 再讨论分支点的连续延拓.

\begin{theorem}[Jost 解与解析列]\label{thm:jost}
设 $q-q_0\in L^1(\mathbb R)\cap L^\infty(\mathbb R)$. 对 $z\in\Sigma\mathbin{\backslash}\{0,\pm\ii q_0\}$, Volterra 方程
\begin{equation}\label{eq:volterra}
\mu_\pm(x,z)=Y(z)+\int_{\pm\infty}^{x}
Y(z)e^{-\ii\lambda(z)(x-y)\sg}Y(z)^{-1}
\Delta Q(y)\mu_\pm(y,z)e^{\ii\lambda(z)(x-y)\sg}\,\dd y
\end{equation}
有唯一关于 $x$ 有界的解. 列 $\mu_{-,1},\mu_{+,2}$ 在 $D_+$ 解析, 列 $\mu_{+,1},\mu_{-,2}$ 在 $D_-$ 解析, 并连续延拓到各自闭域中避开 $0,\pm\ii q_0$ 的紧集. 在共振原像之外, $\phi_\pm=\mu_\pm e^{\ii\theta\sg}$ 是满足 \eqref{eq:jost-normalization} 的唯一空间 Jost 解.
\end{theorem}
\begin{proof}
修正空间方程及其背景恒等式为
\[
(\mu_\pm)_x=U\mu_\pm+\ii\lambda\mu_\pm\sg,\qquad U_\infty Y=-\ii\lambda Y\sg.
\]

\par\medskip\noindent\textbf{把变分常数公式展开成积分因子的计算.}\quad 
在当前的 $z$ 上, $Y$ 与 $x$ 无关且可逆. 把
$U=U_\infty+\Delta Q$ 代入上式, 再左乘 $Y^{-1}$, 得
\begin{align*}
(Y^{-1}\mu_\pm)_x
&=Y^{-1}U_\infty\mu_\pm
+Y^{-1}\Delta Q\mu_\pm
+\ii\lambda Y^{-1}\mu_\pm\sg\\
&=-\ii\lambda\sg Y^{-1}\mu_\pm
+\ii\lambda Y^{-1}\mu_\pm\sg
+Y^{-1}\Delta Q\mu_\pm.
\end{align*}
这里用到的背景恒等式在乘以 $Y^{-1}$ 以后为
\[
Y^{-1}U_\infty Y=-\ii\lambda\sg,
\qquad
Y^{-1}U_\infty=-\ii\lambda\sg Y^{-1}.
\]
现在左右同时乘积分因子. 对整个乘积求导:
\begin{align*}
&\partial_x\bigl(e^{\ii\lambda x\sg}Y^{-1}\mu_\pm
 e^{-\ii\lambda x\sg}\bigr)\\
&=\ii\lambda\sg e^{\ii\lambda x\sg}Y^{-1}\mu_\pm
 e^{-\ii\lambda x\sg}
+e^{\ii\lambda x\sg}(Y^{-1}\mu_\pm)_x
 e^{-\ii\lambda x\sg}\\
&\quad-e^{\ii\lambda x\sg}Y^{-1}\mu_\pm
 \ii\lambda\sg e^{-\ii\lambda x\sg}\\
&=e^{\ii\lambda x\sg}Y^{-1}\Delta Q(x)\mu_\pm(x)
 e^{-\ii\lambda x\sg}.
\end{align*}
第一行求导产生的左侧项, 恰好与微分方程中的
$-\ii\lambda\sg Y^{-1}\mu_\pm$ 相消;
右侧的两项也恰好相消. 从相应的归一化端积分到 $x$, 得
\begin{align*}
&e^{\ii\lambda x\sg}Y^{-1}\mu_\pm(x)e^{-\ii\lambda x\sg}\\
&=I+\int_{\pm\infty}^{x}
 e^{\ii\lambda y\sg}Y^{-1}\Delta Q(y)\mu_\pm(y)
 e^{-\ii\lambda y\sg}\,\dd y.
\end{align*}
再左乘 $Ye^{-\ii\lambda x\sg}$, 右乘
$e^{\ii\lambda x\sg}$, 得
\begin{align*}
\mu_\pm(x)
&=Y+\int_{\pm\infty}^{x}
 Ye^{-\ii\lambda(x-y)\sg}Y^{-1}
 \Delta Q(y)\mu_\pm(y)e^{\ii\lambda(x-y)\sg}\,\dd y.
\end{align*}
对 $+\infty$ 端, 积分方向应原样保留:
\[
\int_{+\infty}^{x}(\cdots)\,\dd y
=-\int_x^{+\infty}(\cdots)\,\dd y.
\]
这个负号在求导时与下限求导的负号抵消.
在连续谱上积分因子有界, 可以从端点归一化直接得到上式.
向非实谱延拓时应逐列使用下文的有界核, 而不是要求整个
$e^{\ii\lambda x\sg}$ 在无穷远有界.

对 $Y^{-1}\mu_\pm$ 使用变分常数公式, 得到 \eqref{eq:volterra}. 令 $s_1=1,s_2=-1$, 第 $j$ 列的传播核为
\begin{equation}\label{eq:columnkernel}
K_j(r,z)=e^{\ii s_j\lambda r}Y e^{-\ii\lambda r\sg}Y^{-1},\qquad r=x-y.
\end{equation}

\par\medskip\noindent\textbf{四种列核的指数从何而来.}\quad 
在矩阵积分式中取第 $j$ 列, 最右边的对角矩阵只给该列乘
$e^{\ii s_j\lambda(x-y)}$. 因此有
\[
\mu_{\pm,j}(x)
=Y_j+\int_{\pm\infty}^{x}
K_j(x-y,z)\Delta Q(y)\mu_{\pm,j}(y)\,\dd y.
\]
把两个指数相乘, 第一列和第二列分别给出
\begin{align*}
K_1(r,z)
&=Y\begin{pmatrix}e^{\ii\lambda r}e^{-\ii\lambda r}&0\\
0&e^{\ii\lambda r}e^{\ii\lambda r}\end{pmatrix}Y^{-1}
=Y\begin{pmatrix}1&0\\0&e^{2\ii\lambda r}\end{pmatrix}Y^{-1},\\
K_2(r,z)
&=Y\begin{pmatrix}e^{-\ii\lambda r}e^{-\ii\lambda r}&0\\
0&e^{-\ii\lambda r}e^{\ii\lambda r}\end{pmatrix}Y^{-1}
=Y\begin{pmatrix}e^{-2\ii\lambda r}&0\\0&1\end{pmatrix}Y^{-1}.
\end{align*}
计算模时要同时看积分方向与谱根的虚部:
\[
|e^{2\ii\lambda r}|=e^{-2r\operatorname{Im}\lambda},
\qquad
|e^{-2\ii\lambda r}|=e^{2r\operatorname{Im}\lambda}.
\]
四种情况逐一为
\begin{align*}
\mu_{-,1}:&\quad y\le x,\ r\ge0,\ \operatorname{Im}\lambda\ge0,
& e^{-2r\operatorname{Im}\lambda}&\le1,\\
\mu_{+,2}:&\quad y\ge x,\ r\le0,\ \operatorname{Im}\lambda\ge0,
& e^{2r\operatorname{Im}\lambda}&\le1,\\
\mu_{+,1}:&\quad y\ge x,\ r\le0,\ \operatorname{Im}\lambda\le0,
& e^{-2r\operatorname{Im}\lambda}&\le1,\\
\mu_{-,2}:&\quad y\le x,\ r\ge0,\ \operatorname{Im}\lambda\le0,
& e^{2r\operatorname{Im}\lambda}&\le1.
\end{align*}
前两列属于 $D_+$, 后两列属于 $D_-$.
例如 $\mu_{+,1}$ 不能沿 $D_+$ 用同一个正端积分式构造,
因为这时 $r\le0$ 而 $\operatorname{Im}\lambda>0$,
$e^{-2r\operatorname{Im}\lambda}$ 在远端增长.
这正是必须区分列而不能给整个 Jost 矩阵指定一个解析域的原因.

以 $e_1,e_2$ 表示标准基向量. 对于负端第一列, 令 $w=Y^{-1}\mu_{-,1}$, 则
\begin{equation}\label{eq:jost-column-volterra}
w(x,z)=e_1+\int_{-\infty}^{x}G(x-y,z)\Delta Q(y)Y(z)w(y,z)\,\dd y,
\end{equation}
其中
\begin{equation}\label{eq:jost-Gkernel}
G(r,z)=\diag(1,e^{2\ii\lambda r})Y^{-1}
=\dfrac{1}{\gamma(z)}
\begin{pmatrix}
1&\dfrac{\ii q_0}{z}\\[4pt]
\dfrac{\ii q_0}{z}e^{2\ii\lambda r}&e^{2\ii\lambda r}
\end{pmatrix}.
\end{equation}
当 $y\le x$, $\operatorname{Im}\lambda\ge0$ 时, $|e^{2\ii\lambda(x-y)}|\le1$, 因而此列对应 $D_+$. 对第二列, 相应对角指数为 $\diag(e^{-2\ii\lambda r},1)$; 它在 $y\ge x$, $\operatorname{Im}\lambda\ge0$ 时有界, 给出 $\mu_{+,2}$ 的解析域. 其余两列取相反的积分方向和谱根虚部.


\par\medskip\noindent\textbf{负端第一列的两个分量积分式.}\quad 
$w=Y^{-1}\mu_{-,1}$ 沿用原文记号. 先把进入积分的矩阵乘法写全:
\[
\Delta Q(y)=\begin{pmatrix}0&q(y)-q_0\\
-\conj{q(y)}+q_0&0\end{pmatrix},
\qquad
Yw=\begin{pmatrix}
w_1-\dfrac{\ii q_0}{z}w_2\\
-\dfrac{\ii q_0}{z}w_1+w_2
\end{pmatrix}.
\]
因此
\[
\Delta Q(y)Yw
=\begin{pmatrix}
(q-q_0)\left(-\dfrac{\ii q_0}{z}w_1+w_2\right)\\
-(\conj{q}-q_0)\left(w_1-\dfrac{\ii q_0}{z}w_2\right)
\end{pmatrix}.
\]
在这里所有 $q,w$ 均取 $y$ 处的值. 再左乘 $G$, 得
\begin{align*}
w_1(x)
&=1+\dfrac{1}{\gamma}\int_{-\infty}^{x}
(q-q_0)\left(-\dfrac{\ii q_0}{z}w_1+w_2\right)\,\dd y\\
&\quad-\dfrac{\ii q_0}{z\gamma}\int_{-\infty}^{x}(\conj{q}-q_0)
\left(w_1-\dfrac{\ii q_0}{z}w_2\right)\,\dd y,\\
w_2(x)
&=\dfrac{\ii q_0}{z\gamma}\int_{-\infty}^{x}e^{2\ii\lambda(x-y)}(q-q_0)
\left(-\dfrac{\ii q_0}{z}w_1+w_2\right)\,\dd y\\
&\quad-\dfrac{1}{\gamma}\int_{-\infty}^{x}e^{2\ii\lambda(x-y)}
(\conj{q}-q_0)\left(w_1-\dfrac{\ii q_0}{z}w_2\right)\,\dd y.
\end{align*}
这两个式子显示: 背景列 $e_1$ 的第二分量为零,
而唯一可能增长的指数只出现在第二分量的传播中.
同样的矩阵乘法对第二列将指数置于第一行.

采用向量范数 $\|v\|_1=|v_1|+|v_2|$ 及其诱导矩阵范数. 由
\[
\|\Delta Q(y)\|_1=|q(y)-q_0|,\qquad
\|Y\|_1=1+\dfrac{q_0}{|z|},\qquad
\|Y^{-1}\|_1=\dfrac{1+\dfrac{q_0}{|z|}}{|\gamma(z)|},
\]
可得 HMB 积分核的具体估计
\begin{equation}\label{eq:jost-kernel-estimate}
\|G(x-y,z)\Delta Q(y)Y\|_1
\le c(z)|q(y)-q_0|,\qquad
c(z)=\dfrac{\left(1+\dfrac{q_0}{|z|}\right)^2}{\left|1+\dfrac{q_0^2}{z^2}\right|}.
\end{equation}
取 $\mathcal K\subset\overline{D_+}\mathbin{\backslash}\{0,\pm\ii q_0\}$ 为紧集, 记
\[
c_{\mathcal K}=\max_{z\in\mathcal K}c(z),\qquad
M_{\mathcal K}(x)=c_{\mathcal K}\int_{-\infty}^{x}|q(y)-q_0|\,\dd y.
\]

\par\medskip\noindent\textbf{矩阵范数估计中的每一个因子.}\quad 
诱导 $1$ 范数是各列绝对值之和的最大值. 由上面三个矩阵的列和,
\begin{align*}
\|\Delta Q\|_1
&=\max\{|\conj{q}-q_0|,|q-q_0|\}=|q-q_0|,\\
\|Y\|_1
&=\max\left\{1+\dfrac{q_0}{|z|},1+\dfrac{q_0}{|z|}\right\}
=1+\dfrac{q_0}{|z|},\\
\|G(r,z)\|_1
&\le\dfrac{1}{|\gamma|}
\max\left\{1+\dfrac{q_0}{|z|},\dfrac{q_0}{|z|}+1\right\}
=\dfrac{1+\dfrac{q_0}{|z|}}{|\gamma|}.
\end{align*}
最后一行使用了允许方向上的 $|e^{2\ii\lambda r}|\le1$.
接着按乘法次序估计,
\begin{align*}
\|G\Delta QYw\|_1
&\le\|G\|_1\|\Delta Q\|_1\|Y\|_1\|w\|_1\\
&\le\dfrac{\left(1+\dfrac{q_0}{|z|}\right)^2}{|\gamma|}
|q-q_0|\|w\|_1.
\end{align*}
紧集避开 $z=0$ 与 $\gamma=0$ 的两点, 因而 $c_{\mathcal K}$ 有限.
$M_{\mathcal K}$ 是非减的绝对连续函数, 且几乎处处满足
\[
M_{\mathcal K}'(x)=c_{\mathcal K}|q(x)-q_0|,
\qquad M_{\mathcal K}(-\infty)=0.
\]
这一个导数等式正是下文阶乘分母的来源.

构造 Neumann 级数
\begin{equation}\label{eq:jost-neumann}
\begin{aligned}
w&=\sum\limits_{n=0}^{\infty}w^{(n)},\qquad w^{(0)}=e_1,\\
w^{(n+1)}&=\int_{-\infty}^{x}G(x-y,z)\Delta Q(y)Yw^{(n)}(y,z)\,\dd y.
\end{aligned}
\end{equation}
若 $\|w^{(n)}(y,z)\|_1\le \dfrac{M_{\mathcal K}(y)^n}{n!}$, 则
\[
\|w^{(n+1)}(x,z)\|_1\le\dfrac{1}{n!}\int_{-\infty}^{x}c_{\mathcal K}|q(y)-q_0|M_{\mathcal K}(y)^n\,\dd y
=\dfrac{M_{\mathcal K}(x)^{n+1}}{(n+1)!}.
\]
归纳即得
\begin{equation}\label{eq:jost-factorial}
\|w^{(n)}(x,z)\|_1\le\dfrac{M_{\mathcal K}(x)^n}{n!},\qquad
\|w(x,z)-e_1\|_1\le e^{M_{\mathcal K}(x)}-1.
\end{equation}

\par\medskip\noindent\textbf{把逐次积分真正展开到阶乘.}\quad 
第一次与第二次迭代分别满足
\begin{align*}
\|w^{(1)}(x,z)\|_1
&\le\int_{-\infty}^{x}M_{\mathcal K}'(y)\,\dd y
=M_{\mathcal K}(x),\\
\|w^{(2)}(x,z)\|_1
&\le\int_{-\infty}^{x}M_{\mathcal K}'(y)
\int_{-\infty}^{y}M_{\mathcal K}'(y_1)\,\dd y_1\,\dd y\\
&=\int_{-\infty}^{x}M_{\mathcal K}'(y)M_{\mathcal K}(y)\,\dd y
=\dfrac{M_{\mathcal K}(x)^2}{2}.
\end{align*}
第三次同样得到
\[
\|w^{(3)}(x,z)\|_1
\le\int_{-\infty}^{x}M_{\mathcal K}'(y)
\dfrac{M_{\mathcal K}(y)^2}{2}\,\dd y
=\dfrac{M_{\mathcal K}(x)^3}{6}.
\]
一般一步的等号由微积分基本定理逐行给出:
\begin{align*}
\dfrac{1}{n!}\int_{-\infty}^{x}
M_{\mathcal K}'(y)M_{\mathcal K}(y)^n\,\dd y
&=\dfrac{1}{n!}\int_{-\infty}^{x}
\dfrac{1}{n+1}\partial_y\bigl(M_{\mathcal K}(y)^{n+1}\bigr)\,\dd y\\
&=\dfrac{M_{\mathcal K}(x)^{n+1}-M_{\mathcal K}(-\infty)^{n+1}}{(n+1)!}\\
&=\dfrac{M_{\mathcal K}(x)^{n+1}}{(n+1)!}.
\end{align*}
等价地, 第 $n$ 次迭代的积分区域是
\[
-\infty<y_n<y_{n-1}<\cdots<y_1<x.
\]
把 $( -\infty,x)^n$ 按 $n$ 个坐标的先后次序分成 $n!$ 个区域;
被积的范数上界是各坐标中同一非负函数的乘积, 因而每个区域的积分相等.
于是其积分为整个乘积积分的 $\dfrac{1}{n!}$, 即
$\dfrac{M_{\mathcal K}(x)^n}{n!}$.
求和时从 $n=1$ 开始的部分为
\[
\sum\limits_{n=1}^{\infty}\dfrac{M_{\mathcal K}(x)^n}{n!}
=e^{M_{\mathcal K}(x)}-1.
\]
特别地, 不要求 $\|q-q_0\|_{L^1}$ 小于某个数;
阶乘使任意有限的 $L^1$ 质量都能给出收敛.

由于 $M_{\mathcal K}(x)\le c_{\mathcal K}\|q-q_0\|_{L^1}$, 级数对实 $x$ 及 $z\in\mathcal K$ 绝对一致收敛, 且在 $x\to-\infty$ 时趋于 $e_1$. 每一项的积分核在 $D_+$ 解析, 并受上述可积函数控制; 参数积分定理及紧集一致收敛给出 $w$ 和 $\mu_{-,1}=Yw$ 的解析性. 同一控制函数与支配收敛定理给出闭域连续性.


\par\medskip\noindent\textbf{参数解析性与闭域连续性所用的控制.}\quad 
对任一固定有限的 $x$, 先把积分在 $y=-R$ 处截断.
在有限区间上, 核的各矩阵元都是 $z$ 的解析函数,
因为 $Y,Y^{-1},\lambda$ 在当前紧集附近解析, 指数是整函数.
再让 $R\to+\infty$, 积分尾部的范数至多为
\[
c_{\mathcal K}\sup\limits_{y,z}\|w^{(n)}(y,z)\|_1
\int_{-\infty}^{-R}|q(y)-q_0|\,\dd y,
\]
它与 $z\in\mathcal K$ 无关并趋于零.
所以每一项 $w^{(n)}$ 都是解析函数的一致极限.
对所有迭代项再使用
\[
\sum\limits_{n=0}^{\infty}
\sup\limits_{x\in\mathbb R,z\in\mathcal K}\|w^{(n)}(x,z)\|_1
\le\sum\limits_{n=0}^{\infty}
\dfrac{\bigl(c_{\mathcal K}\|q-q_0\|_{L^1}\bigr)^n}{n!}<\infty.
\]
解析函数列的紧集一致极限仍然解析.
若 $z$ 从解析域趋向闭域边界, 核逐点连续,
上面同一个可积上界仍适用; 对每项与级数分别取极限,
即可交换边界极限、积分和无穷求和.
这里不声称仅靠 $L^1$ 就能在边界上任意次求 $z$ 导数;
对导数直接估计会出现 $|x-y|$, 需要与所求结论相适应的矩条件.

对两个有界解之差 $d$ 迭代 $n$ 次, 有
\[
\|d(x,z)\|_1\le\|d(\cdot,z)\|_{L^\infty}\dfrac{M_{\mathcal K}(x)^n}{n!}\longrightarrow0,
\]
从而解唯一. 
\par\medskip\noindent\textbf{唯一性与空间微分各自怎样完成.}\quad 
差 $d$ 的自由项为零. 第一次应用积分方程给出
\[
\|d(x,z)\|_1\le\|d\|_{L^\infty}M_{\mathcal K}(x).
\]
把同一个积分式再次代入积分中的 $d(y,z)$,
积分区域成为 $y_1\le x,y_2\le y_1$, 故
\[
\|d(x,z)\|_1\le\|d\|_{L^\infty}
\dfrac{M_{\mathcal K}(x)^2}{2!}.
\]
重复 $n$ 次后就是原文所列估计. 对固定 $M_{\mathcal K}(x)$,
相邻上界之比为 $\dfrac{M_{\mathcal K}(x)}{n+1}$, 趋于零,
因此右端趋于零, 从而 $d(x,z)=0$.
对空间导数, 直接使用原文的列核有
\[
K_j(0,z)=I,
\qquad
\partial_rK_j(r,z)=(U_\infty+\ii s_j\lambda I)K_j(r,z).
\]
这是因为 $Ye^{-\ii\lambda r\sg}Y^{-1}=e^{rU_\infty}$.
逐列对积分式求导, 可得
\begin{align*}
(\mu_{\pm,j})_x
&=\Delta Q(x)\mu_{\pm,j}(x)
+\int_{\pm\infty}^{x}(U_\infty+\ii s_j\lambda I)
 K_j(x-y,z)\Delta Q(y)\mu_{\pm,j}(y)\,\dd y\\
&=\Delta Q\mu_{\pm,j}
+(U_\infty+\ii s_j\lambda I)(\mu_{\pm,j}-Y_j)\\
&=(U+\ii s_j\lambda I)\mu_{\pm,j}.
\end{align*}
最后使用 $(U_\infty+\ii s_j\lambda I)Y_j=0$.
故两列合起来满足修正空间方程.
再计算恢复指数后的导数:
\begin{align*}
\partial_x(\mu_{\pm,j}e^{\ii s_j\theta})
&=(U\mu_{\pm,j}+\ii s_j\lambda\mu_{\pm,j})e^{\ii s_j\theta}
-\ii s_j\lambda\mu_{\pm,j}e^{\ii s_j\theta}\\
&=U\phi_{\pm,j}.
\end{align*}
在目前的低正则性假设下, 上述空间导数理解为几乎处处导数,
解在有限区间上绝对连续; 若 $q$ 连续, 等式处处成立.

正端取 $\displaystyle c_{\mathcal K}\int_x^{+\infty}|q(y)-q_0|\,\dd y$, 反向有序积分给出相同的阶乘估计; 另外两列亦同. 在 $\Sigma$ 上四列均可取边界值, 构成 \eqref{eq:volterra} 的有界矩阵解. 对积分方程作局部绝对连续意义下的 $x$ 微分, 恢复空间谱方程, 其归一化由 \eqref{eq:jost-factorial} 给出. 任一具有同一振荡归一化的空间解通过变分常数公式满足该 Volterra 方程, 因而也唯一.
\end{proof}

\begin{proposition}[向分支点的连续延拓]\label{prop:jost-branch}
在定理 \ref{thm:jost} 的条件上, 再设
\begin{equation}\label{eq:jost-weighted}
\int_{-\infty}^{+\infty}(1+|x|)|q(x)-q_0|\,\dd x<\infty.
\end{equation}
则每个 Jost 解析列均可从其闭解析域连续延拓到 $z_b=\pm\ii q_0$. 端点值满足
\begin{equation}\label{eq:jost-branch-volterra}
\mu_\pm(x,z_b)=Y(z_b)+\int_{\pm\infty}^{x}
[I+(x-y)U_\infty(z_b)]\Delta Q(y)\mu_\pm(y,z_b)\,\dd y,
\end{equation}
并在对应归一化端趋于 $Y(z_b)$. 此矩阵在任一对应半直线上有界, 满足空间谱方程, 且秩为一.
\end{proposition}
\begin{proof}

\par\medskip\noindent\textbf{为什么分支点必须重新估计核.}\quad 
当 $z\to z_b$, $\gamma(z)\to0$, 因而前面的
$c(z)$ 发散. 这只说明将 $Y$ 与 $Y^{-1}$ 分开估计过粗,
并不说明二者夹住指数后的整体发散.
先把背景矩阵的平方完整相乘:
\[
\begin{pmatrix}-\ii k&q_0\\-q_0&\ii k\end{pmatrix}^2
=\begin{pmatrix}
-k^2-q_0^2&-\ii kq_0+\ii kq_0\\
\ii kq_0-\ii kq_0&-q_0^2-k^2
\end{pmatrix}
=-\lambda^2I.
\]
所以偶数次幂与奇数次幂分别为
\[
U_\infty^{2n}=(-1)^n\lambda^{2n}I,
\qquad
U_\infty^{2n+1}=(-1)^n\lambda^{2n}U_\infty.
\]
代入矩阵指数的幂级数并把偶奇两部分分开,
\begin{align*}
e^{rU_\infty}
&=\sum\limits_{n=0}^{\infty}\dfrac{(-1)^n(\lambda r)^{2n}}{(2n)!}I
+\sum\limits_{n=0}^{\infty}\dfrac{(-1)^n\lambda^{2n}r^{2n+1}}{(2n+1)!}U_\infty\\
&=\cos(\lambda r)I+\dfrac{\sin(\lambda r)}{\lambda}U_\infty.
\end{align*}
在 $\lambda=0$ 处第二个系数按连续值取 $r$.

由 $U_\infty^2=-\lambda^2I$, 背景传播矩阵可写为
\begin{equation}\label{eq:background-propagator}
e^{rU_\infty}
=\cos(\lambda r)I+\dfrac{\sin(\lambda r)}{\lambda}U_\infty
\longrightarrow I+rU_\infty(z_b),\qquad z\to z_b.
\end{equation}
尽管 $Y^{-1}$ 在分支点发散, \eqref{eq:columnkernel} 中的乘积具有可控的连续极限. 具体地,
\begin{equation}\label{eq:branch-columnkernel}
K_j(r,z)=\dfrac{1+e^{2\ii s_j\lambda r}}{2}I
+\dfrac{e^{2\ii s_j\lambda r}-1}{2\ii s_j\lambda}U_\infty(z).
\end{equation}

\par\medskip\noindent\textbf{消掉分支点的表面分母.}\quad 
将上式乘上列核中的 $e^{\ii s_j\lambda r}$.
利用 $s_j=\pm1$, 分别计算
\begin{align*}
e^{\ii s_j\lambda r}\cos(\lambda r)
&=\dfrac{1+e^{2\ii s_j\lambda r}}{2},\\
e^{\ii s_j\lambda r}\dfrac{\sin(\lambda r)}{\lambda}
&=\dfrac{e^{2\ii s_j\lambda r}-1}{2\ii s_j\lambda}.
\end{align*}
这得到原文的列核表达式. 第二个系数不是实质的极点:
\begin{align*}
e^{2\ii s_j\lambda r}-1
&=\int_0^1\partial_\tau
 e^{2\ii s_j\tau\lambda r}\,\dd\tau\\
&=2\ii s_j\lambda r\int_0^1
 e^{2\ii s_j\tau\lambda r}\,\dd\tau.
\end{align*}
当 $\lambda\to0$ 时, 这个积分趋于 1, 因而商趋于 $r$.
如果 $s_j\operatorname{Im}\lambda\,r\ge0$, 则每个
$0\le\tau\le1$ 上的指数模均不超过 1, 于是
\[
\left|\dfrac{1+e^{2\ii s_j\lambda r}}{2}\right|\le1,
\qquad
\left|\dfrac{e^{2\ii s_j\lambda r}-1}{2\ii s_j\lambda}\right|
\le|r|.
\]
故只需控制 $U_\infty$ 在分支点邻域的范数,
就能得到 $1+C_b|r|$ 的一致上界.
这解释了一阶空间矩为何足以处理分支点.

在该列允许的积分方向及闭解析域内, $s_j\operatorname{Im}\lambda\,r\ge0$, 且
\[
\dfrac{e^{2\ii s_j\lambda r}-1}{2\ii s_j\lambda}
=r\int_0^1 e^{2\ii s_j\tau\lambda r}\,\dd\tau,\qquad
\|K_j(r,z)\|_1\le1+C_b|r|.
\]
这里 $C_b$ 可在 $z_b$ 的一个闭邻域上一致选取. 这同时给出 $K_j(r,z_b)=I+rU_\infty(z_b)$.

以负端为例, 固定 $b\in\mathbb R$, 对 $y\le x\le b$ 有
\[
1+C_b|x-y|\le A_b(1+|y|),\qquad
L_b(x)=A_b\int_{-\infty}^{x}(1+|y|)|q(y)-q_0|\,\dd y.
\]

\par\medskip\noindent\textbf{加权估计与半直线归一化.}\quad 
当 $y\le x\le b$ 时,
\[
0\le x-y\le b-y\le|b|+|y|,
\]
故
\[
1+C_b|x-y|
\le 1+C_b|b|+C_b|y|
\le (1+C_b|b|+C_b)(1+|y|).
\]
因而可以把 $A_b$ 选为括号中的常数或更大的常数.
这不是对整个实轴统一控制 $|x-y|$;
它是在任意固定半直线 $x\le b$ 上控制积分核.
对第 $n$ 次迭代应用与前面完全相同的有序积分,
每个 $\Delta Q$ 的范数与 $1+C_b|x-y|$ 相乘后,
都由 $A_b(1+|y|)|q(y)-q_0|$ 控制.
于是阶乘估计的积分函数变为 $L_b$, 不再出现 $\dfrac{1}{\gamma}$.
又有
\[
\|\mu_{-,j}(x,z)-Y_j(z)\|_1
\le\left(\sup\limits_z\|Y_j(z)\|_1\right)(e^{L_b(x)}-1).
\]
加权可积性给出 $L_b(x)\to0$ 当 $x\to-\infty$,
所以端点归一化与 $z\to z_b$ 极限可以兼容.
在正端, 对 $b\le x\le y$ 用
$y-x\le |y|+|b|$, 并把积分上界改为从 $x$ 到 $+\infty$,
其余每一步保持相同.

直接对 $\mu_{-,j}$ 的列积分方程迭代, 其第 $n$ 项满足
\[
\|\mu_{-,j}^{(n)}(x,z)\|_1
\le\left(\sup_{z}\|Y_j(z)\|_1\right)\dfrac{L_b(x)^n}{n!},\qquad x\le b.
\]
加权可积条件保证级数在 $x\le b$ 及分支点邻域内一致收敛. 各项利用 \eqref{eq:branch-columnkernel} 支配收敛, 故其和连续趋于 \eqref{eq:jost-branch-volterra}; $L_b(x)\to0$ 则给出负端归一化. 正端使用 $y\ge x\ge b$ 和从 $x$ 到 $+\infty$ 的加权积分, 得到相同结论. 对差的有序积分估计还给出端点 Volterra 解在对应半直线上的唯一性.


\par\medskip\noindent\textbf{端点空间方程与秩一结论的具体计算.}\quad 
在 $z_b=\ii q_0$ 处,
\[
Y(z_b)=\begin{pmatrix}1&-1\\-1&1\end{pmatrix},
\qquad Y_2(z_b)=-Y_1(z_b).
\]
在 $z_b=-\ii q_0$ 处,
\[
Y(z_b)=\begin{pmatrix}1&1\\1&1\end{pmatrix},
\qquad Y_2(z_b)=Y_1(z_b).
\]
这两种线性关系都能逐列代入同一个端点 Volterra 方程.
负号情形下, $\mu_{\pm,2}+\mu_{\pm,1}$ 的自由项为零;
正号情形下, $\mu_{\pm,2}-\mu_{\pm,1}$ 的自由项为零.
由该半直线上的唯一性, 相应组合恒为零.
求端点方程的空间导数时,
\begin{align*}
(\mu_\pm)_x
&=\Delta Q(x)\mu_\pm(x)
+U_\infty\int_{\pm\infty}^{x}
\Delta Q(y)\mu_\pm(y)\,\dd y,\\
U_\infty\mu_\pm(x)
&=U_\infty Y
+\int_{\pm\infty}^{x}
\bigl(U_\infty+(x-y)U_\infty^2\bigr)
\Delta Q(y)\mu_\pm(y)\,\dd y\\
&=U_\infty\int_{\pm\infty}^{x}\Delta Q(y)\mu_\pm(y)\,\dd y.
\end{align*}
因此 $(\mu_\pm)_x=(U_\infty+\Delta Q)\mu_\pm=U\mu_\pm$.
第一列在归一化端趋于上面非零的 $Y_1$.
若它在某个有限点为零, 一阶线性初值问题的唯一性会使它处处为零,
与非零端态矛盾. 所以两列始终相关且至少一列非零, 矩阵秩恰为一.

最后, $U_\infty(z_b)Y(z_b)=0$, $U_\infty(z_b)^2=0$. 对 \eqref{eq:jost-branch-volterra} 求导得到 $(\mu_\pm)_x=U(z_b)\mu_\pm$. 因 $Y(z_b)$ 的两列线性相关, 唯一性使同一关系对 $\mu_\pm(x,z_b)$ 成立. 其归一化端值非零, 而一阶线性系统的非零解不能在有限点为零, 因而矩阵秩恒为一.
\end{proof}

共振原像 $z_a^\pm$ 不使空间积分核退化, 因而 $\mu_\pm$ 在此仍由 \eqref{eq:volterra} 正则取值. 传播相位在共振处的处理及两组 Jost 解的共同演化见第 \ref{ch:scattering} 节.
